\section{De John a Sandy: cada paso del flujo de trabajo}
\subsection{Un ticket ordinario: la larga cadena del proceso manual}
En la historia de 00:02:11 a 00:11:21, John derrama café sobre su laptop por la mañana y abre un ticket con TI para solicitar un reemplazo del equipo; el ticket se asigna a Sandy, quien revisa uno por uno la base de conocimiento y los incidentes históricos, ajusta los permisos y deja un registro claro. Aunque ya está desplegada la IA generativa más reciente, Sandy sigue realizando una gran cantidad de pasos manuales, entre ellos verificar el control de acceso, invocar los sistemas relacionados y escribir el resumen de la solución.

\begin{warningbox}{La brecha del «dominó» de la automatización}
Aunque exista un resumen generado por IA, en los flujos de trabajo empresariales todavía hay cientos de escenarios que requieren criterio humano, verificación de permisos, evaluación de cumplimiento y coordinación entre varios sistemas; cualquier automatización debe conservar en este proceso la trazabilidad y un margen para la intervención humana.
\end{warningbox}

\subsection{Tareas «grano de arena»: poco frecuentes pero cruciales}
El presentador señala que la mayoría de las tareas de «bajo valor, baja frecuencia» son demasiado singulares para que los scripts tradicionales las cubran; estas «tareas grano de arena» se repiten millones de veces al día en las empresas de todo el mundo, pero cada una exige una respuesta rápida.

\begin{importantbox}{El alcance del Agente no debe medirse solo por la «alta frecuencia»
}
El Agente debe aprovechar las oportunidades de ser «explicable, trazable y ajustable con granularidad fina»: aunque una sola tarea tome poco tiempo, lo que realmente importa es el incremento de eficiencia y la satisfacción de los empleados que se generan a escala aggregate.
\end{importantbox}

\subsection{Resumen de la sección}
El ejemplo de John y Sandy muestra el valor potencial y los riesgos del Agente en escenarios con múltiples sistemas, roles y pasos, y sienta las bases para los requisitos de arquitectura y auditoría que se tratan más adelante.
